\begin{abstract}\noindent Token-level credit assignment for language-model
reinforcement learning is an active area of research, with methods that
reweight policy-gradient credit using intrinsic signals such as surprisal,
entropy, and divergence from a reference policy. At the same time, practical
LLM-RL pipelines overwhelmingly rely on parameter-efficient fine-tuning,
especially LoRA. These threads are usually studied separately: token-credit
methods are specified as if the policy were fully trainable, while PEFT is
treated as an implementation detail. We argue that this separation is
misleading. Under LoRA, the policy is restricted to a low-rank neighbourhood of
the reference, which can flatten the per-token differences used by common
intrinsic credit signals and drive their normalized weights toward uniform
broadcast. We formalize this failure mode and propose measuring it directly
with concentration diagnostics such as weight Gini and effective-token ratio.
We then introduce \emph{Adapter-Residual Credit Assignment} (ARCA), a
lightweight alternative that derives token salience from the adapter's own
hidden-state residual,
$\|h^{\text{adapted}}_t - h^{\text{base}}_t\|_2$. This signal asks where the
adapter actually changes the model, rather than where the output distribution
appears uncertain or shifted, and requires no reward model, value head, or tree
construction. Our experiments are designed to test whether adapter-residual
credit remains non-degenerate under LoRA and whether preserving that structure
improves RL fine-tuning under matched rollout budgets.
\end{abstract}
